\documentclass[12pt,a4paper]{article}
\usepackage[spanish]{babel}
\usepackage[utf8]{inputenc}
\usepackage{amsmath,amssymb}
\usepackage[margin=2.2cm]{geometry}
\usepackage{enumitem}
\usepackage{tikz}
\usepackage{capt-of}
\usepackage{parskip}

\usetikzlibrary{arrows.meta,calc}

% Comando corregido para la circunferencia unitaria
\newcommand{\unitcircle}[1]{%
\begin{center}
\begin{tikzpicture}[scale=0.75]
  \draw[->] (-1.5,0) -- (1.5,0) node[right] {$x$};
  \draw[->] (0,-1.5) -- (0,1.5) node[above] {$y$};
  \draw (0,0) circle (1.2);
  \foreach \coord/\lab in {#1} {
    \filldraw (\coord) circle (2pt) node[font=\footnotesize, fill=white, inner sep=1pt] {$\lab$};
  }
\end{tikzpicture}
\captionof{figure}{Soluciones en la circunferencia unitaria en $[0,2\pi)$.}
\end{center}
}

\setlength{\parskip}{4pt plus2pt minus1pt}
\setlength{\parindent}{0pt}

\title{\textbf{Soluciones paso a paso con gráficos}}
\author{}
\date{}

\begin{document}
\maketitle

% Solución 1
\begin{minipage}{\textwidth}
\section*{Ejercicio 1}
\textbf{Ecuación:} \(2\cos^2 x - \cos x - 1 = 0\) en \([0,2\pi)\).

\textbf{Solución:}
\begin{enumerate}[label=\arabic*.]
    \item \textbf{Sustitución:} Sea \(u = \cos x\). Entonces:
    \[
    2u^2 - u - 1 = 0.
    \]
    \item \textbf{Factorización:}
    \[
    2u^2 - u - 1 = (2u + 1)(u - 1) = 0.
    \]
    De donde \(2u + 1 = 0\) o \(u - 1 = 0\).
    \item \textbf{Soluciones para \(u\):}
    \[
    u = -\frac{1}{2} \quad\text{o}\quad u = 1.
    \]
    \item \textbf{Volver a \(\cos x\):}
    \begin{itemize}
        \item \(\cos x = 1\).
        \item \(\cos x = -\frac{1}{2}\).
    \end{itemize}
    \item \textbf{Resolver \(\cos x = 1\):}
    \[
    x = 0 + 2k\pi \quad (k\in\mathbb{Z}).
    \]
    En \([0,2\pi)\) la única solución es \(x = 0\).
    
    \item \textbf{Resolver \(\cos x = -\frac{1}{2}\):}
    Los ángulos de referencia son \(\frac{\pi}{3}\). Como el coseno es negativo en los cuadrantes II y III:
    \begin{align*}
        x &= \pi - \frac{\pi}{3} + 2k\pi = \frac{2\pi}{3} + 2k\pi,\\
        x &= \pi + \frac{\pi}{3} + 2k\pi = \frac{4\pi}{3} + 2k\pi.
    \end{align*}
    En \([0,2\pi)\): \(x = \frac{2\pi}{3}\) y \(x = \frac{4\pi}{3}\).

    \item \textbf{Conjunto solución:}
    \[
    \boxed{0,\ \frac{2\pi}{3},\ \frac{4\pi}{3}}
    \]
\end{enumerate}

\unitcircle{%
  {(1.2,0)/0},%
  {{1.2*cos(120)},{1.2*sin(120)}}/{\frac{2\pi}{3}},%
  {{1.2*cos(240)},{1.2*sin(240)}}/{\frac{4\pi}{3}}%
}
\end{minipage}

\vspace{1cm}

% Solución 2
\begin{minipage}{\textwidth}
\section*{Ejercicio 2}
\textbf{Ecuación:} \(3\tan^2 x - 4\tan x + 1 = 0\) en \([0,2\pi)\).

\textbf{Solución:}
\begin{enumerate}[label=\arabic*.]
    \item \textbf{Sustitución:} \(u = \tan x\) \(\Rightarrow\) \(3u^2 - 4u + 1 = 0\).
    \item \textbf{Factorización:} \((3u - 1)(u - 1) = 0\).
    \item \textbf{Raíces:} \(u = \frac{1}{3}\) o \(u = 1\).
    \item \textbf{Volver a \(\tan x\):}
    \begin{itemize}
        \item \(\tan x = 1\).
        \item \(\tan x = \frac{1}{3}\).
    \end{itemize}
    \item \textbf{Para \(\tan x = 1\):} Solución general: \(x = \frac{\pi}{4} + k\pi\). En \([0,2\pi)\): \(x = \frac{\pi}{4},\ \frac{5\pi}{4}\).
    \item \textbf{Para \(\tan x = \frac{1}{3}\):} No es un ángulo especial. Usamos la función inversa:
    \[
    x = \arctan\!\left(\frac{1}{3}\right) + k\pi.
    \]
    En \([0,2\pi)\) con \(k=0\): \(x_1 = \arctan\!\left(\frac{1}{3}\right)\) (cuadrante I).\\
    Con \(k=1\): \(x_2 = \arctan\!\left(\frac{1}{3}\right) + \pi\) (cuadrante III).
    \item \textbf{Conjunto solución:}
    \[
    \boxed{\frac{\pi}{4},\ \frac{5\pi}{4},\ \arctan\!\left(\frac{1}{3}\right),\ \pi + \arctan\!\left(\frac{1}{3}\right)}
    \]
\end{enumerate}

\unitcircle{%
  {{1.2*cos(45)},{1.2*sin(45)}}/{\frac{\pi}{4}},%
  {{1.2*cos(225)},{1.2*sin(225)}}/{\frac{5\pi}{4}},%
  {{1.2*cos(18.4349)},{1.2*sin(18.4349)}}/{\alpha},%
  {{1.2*cos(198.4349)},{1.2*sin(198.4349)}}/{\pi{+}\alpha}%
}
{\small $\alpha = \arctan\frac13$.}
\end{minipage}

\vspace{1cm}

% Solución 3
\begin{minipage}{\textwidth}
\section*{Ejercicio 3}
\textbf{Ecuación:} \(\sin^2 x + \sin x - 2 = 0\) en \([0,2\pi)\).

\textbf{Solución:}
\begin{enumerate}[label=\arabic*.]
    \item \textbf{Sustitución:} \(u = \sin x\) \(\Rightarrow\) \(u^2 + u - 2 = 0\).
    \item \textbf{Factorización:} \((u + 2)(u - 1) = 0\).
    \item \textbf{Raíces:} \(u = -2\) o \(u = 1\).
    \item \textbf{Análisis:} \(\sin x = -2\) no tiene solución (\(-1 \le \sin x \le 1\)).\\
    \(\sin x = 1\) sí es válida.
    \item \textbf{Resolver \(\sin x = 1\):}
    \[
    x = \frac{\pi}{2} + 2k\pi.
    \]
    En \([0,2\pi)\): \(x = \frac{\pi}{2}\).
    \item \textbf{Conjunto solución:}
    \[
    \boxed{\frac{\pi}{2}}
    \]
\end{enumerate}

\unitcircle{%
  {(0,1.2)/{\frac{\pi}{2}}}%
}
\end{minipage}

\vspace{1cm}

% Solución 4
\begin{minipage}{\textwidth}
\section*{Ejercicio 4}
\textbf{Ecuación:} \(2\cos^2 x - 5\cos x + 2 = 0\) en \([0,2\pi)\).

\textbf{Solución:}
\begin{enumerate}[label=\arabic*.]
    \item \textbf{Sustitución:} \(u = \cos x\) \(\Rightarrow\) \(2u^2 - 5u + 2 = 0\).
    \item \textbf{Factorización:} \((2u - 1)(u - 2) = 0\).
    \item \textbf{Raíces:} \(u = \frac{1}{2}\) o \(u = 2\).
    \item \textbf{Análisis:} \(\cos x = 2\) no tiene solución.\\
    \(\cos x = \frac{1}{2}\) es válida.
    \item \textbf{Resolver \(\cos x = \frac{1}{2}\):}
    Ángulo de referencia \(\frac{\pi}{3}\); coseno positivo en I y IV:
    \begin{align*}
        x &= \frac{\pi}{3} + 2k\pi,\\
        x &= -\frac{\pi}{3} + 2k\pi \quad (\text{equivalentemente } \frac{5\pi}{3} + 2k\pi).
    \end{align*}
    En \([0,2\pi)\): \(x = \frac{\pi}{3},\ \frac{5\pi}{3}\).
    \item \textbf{Conjunto solución:}
    \[
    \boxed{\frac{\pi}{3},\ \frac{5\pi}{3}}
    \]
\end{enumerate}

\unitcircle{%
  ({1.2*cos(60)},{1.2*sin(60)})/\frac{\pi}{3},
  ({1.2*cos(300)},{1.2*sin(300)})/\frac{5\pi}{3}%
}
\end{minipage}

\vspace{1cm}

% Solución 5
\begin{minipage}{\textwidth}
\section*{Ejercicio 5}
\textbf{Ecuación:} \(4\sin^2 x - 4\sin x + 1 = 0\) en \([0,2\pi)\).

\textbf{Solución:}
\begin{enumerate}[label=\arabic*.]
    \item \textbf{Sustitución:} \(u = \sin x\) \(\Rightarrow\) \(4u^2 - 4u + 1 = 0\).
    \item \textbf{Factorización:} \((2u - 1)^2 = 0\).
    \item \textbf{Raíz doble:} \(u = \frac{1}{2}\).
    \item \textbf{Resolver \(\sin x = \frac{1}{2}\):}
    Ángulo de referencia \(\frac{\pi}{6}\); seno positivo en I y II:
    \begin{align*}
        x &= \frac{\pi}{6} + 2k\pi,\\
        x &= \pi - \frac{\pi}{6} + 2k\pi = \frac{5\pi}{6} + 2k\pi.
    \end{align*}
    En \([0,2\pi)\): \(x = \frac{\pi}{6},\ \frac{5\pi}{6}\).
    \item \textbf{Conjunto solución:}
    \[
    \boxed{\frac{\pi}{6},\ \frac{5\pi}{6}}
    \]
\end{enumerate}

\unitcircle{%
  ({1.2*cos(30)},{1.2*sin(30)})/\frac{\pi}{6},
  ({1.2*cos(150)},{1.2*sin(150)})/\frac{5\pi}{6}%
}
\end{minipage}

\vspace{1cm}

% Solución 6
\begin{minipage}{\textwidth}
\section*{Ejercicio 6}
\textbf{Ecuación:} \(\sec^2 x - 3\sec x + 2 = 0\) en \([0,2\pi)\).

\textbf{Solución:}
\begin{enumerate}[label=\arabic*.]
    \item \textbf{Sustitución:} \(u = \sec x\) \(\Rightarrow\) \(u^2 - 3u + 2 = 0\).
    \item \textbf{Factorización:} \((u - 1)(u - 2) = 0\).
    \item \textbf{Raíces:} \(u = 1\) o \(u = 2\).
    \item \textbf{Volver a \(\cos x\):} Como \(\sec x = \frac{1}{\cos x}\),
    \begin{itemize}
        \item \(\sec x = 1 \Rightarrow \cos x = 1\).
        \item \(\sec x = 2 \Rightarrow \cos x = \frac{1}{2}\).
    \end{itemize}
    \item \textbf{Resolver \(\cos x = 1\):}
    \(x = 0 + 2k\pi\). En \([0,2\pi)\): \(x = 0\).
    \item \textbf{Resolver \(\cos x = \frac{1}{2}\):}
    \begin{align*}
        x &= \frac{\pi}{3} + 2k\pi,\\
        x &= \frac{5\pi}{3} + 2k\pi.
    \end{align*}
    En \([0,2\pi)\): \(x = \frac{\pi}{3},\ \frac{5\pi}{3}\).
    \item \textbf{Conjunto solución:}
    \[
    \boxed{0,\ \frac{\pi}{3},\ \frac{5\pi}{3}}
    \]
\end{enumerate}

\unitcircle{%
  {(1.2,0)/0},
  ({1.2*cos(60)},{1.2*sin(60)})/\frac{\pi}{3},
  ({1.2*cos(300)},{1.2*sin(300)})/\frac{5\pi}{3}%
}
\end{minipage}

\vspace{1cm}

% Solución 7
\begin{minipage}{\textwidth}
\section*{Ejercicio 7}
\textbf{Ecuación:} \(2\csc^2 x - 3\csc x - 2 = 0\) en \([0,2\pi)\).

\textbf{Solución:}
\begin{enumerate}[label=\arabic*.]
    \item \textbf{Sustitución:} \(u = \csc x\) \(\Rightarrow\) \(2u^2 - 3u - 2 = 0\).
    \item \textbf{Factorización:} \((2u + 1)(u - 2) = 0\).
    \item \textbf{Raíces:} \(u = -\frac{1}{2}\) o \(u = 2\).
    \item \textbf{Volver a \(\sin x\):} \(\csc x = \frac{1}{\sin x}\).
    \begin{itemize}
        \item \(\csc x = -\frac{1}{2} \Rightarrow \sin x = -2\) (no válido, \(|\sin x|\le 1\)).
        \item \(\csc x = 2 \Rightarrow \sin x = \frac{1}{2}\).
    \end{itemize}
    \item \textbf{Resolver \(\sin x = \frac{1}{2}\):}
    \begin{align*}
        x &= \frac{\pi}{6} + 2k\pi,\\
        x &= \frac{5\pi}{6} + 2k\pi.
    \end{align*}
    En \([0,2\pi)\): \(x = \frac{\pi}{6},\ \frac{5\pi}{6}\).
    \item \textbf{Conjunto solución:}
    \[
    \boxed{\frac{\pi}{6},\ \frac{5\pi}{6}}
    \]
\end{enumerate}

\unitcircle{%
  ({1.2*cos(30)},{1.2*sin(30)})/\frac{\pi}{6},
  ({1.2*cos(150)},{1.2*sin(150)})/\frac{5\pi}{6}%
}
\end{minipage}

\vspace{1cm}

% Solución 8
\begin{minipage}{\textwidth}
\section*{Ejercicio 8}
\textbf{Ecuación:} \(3\sin^2 x + \sin x - 2 = 0\) en \([0,2\pi)\).

\textbf{Solución:}
\begin{enumerate}[label=\arabic*.]
    \item \textbf{Sustitución:} \(u = \sin x\) \(\Rightarrow\) \(3u^2 + u - 2 = 0\).
    \item \textbf{Factorización:}
    \[
    3u^2 + u - 2 = (3u - 2)(u + 1) = 0.
    \]
    \item \textbf{Raíces:} \(u = \frac{2}{3}\) o \(u = -1\).
    \item \textbf{Resolver \(\sin x = -1\):}
    \[
    x = \frac{3\pi}{2} + 2k\pi.
    \]
    En \([0,2\pi)\): \(x = \frac{3\pi}{2}\).
    \item \textbf{Resolver \(\sin x = \frac{2}{3}\):} No es ángulo especial. Usamos la función \(\arcsin\):
    \begin{align*}
        x &= \arcsin\!\left(\frac{2}{3}\right) + 2k\pi,\\
        x &= \pi - \arcsin\!\left(\frac{2}{3}\right) + 2k\pi.
    \end{align*}
    En \([0,2\pi)\) con \(k=0\):
    \[
    x_1 = \arcsin\!\left(\frac{2}{3}\right) \quad\text{(cuadrante I)},\qquad
    x_2 = \pi - \arcsin\!\left(\frac{2}{3}\right) \quad\text{(cuadrante II)}.
    \]
    \item \textbf{Conjunto solución:}
    \[
    \boxed{\frac{3\pi}{2},\ \arcsin\!\left(\frac{2}{3}\right),\ \pi - \arcsin\!\left(\frac{2}{3}\right)}
    \]
\end{enumerate}

\unitcircle{%
  {(0,-1.2)/\frac{3\pi}{2}},
  ({1.2*cos(41.8103)},{1.2*sin(41.8103)})/\beta,
  ({1.2*cos(138.1897)},{1.2*sin(138.1897)})/\pi{-}\beta%
} \\[2mm]
{\small $\beta = \arcsin\frac23$.}
\end{minipage}

\vspace{1cm}

% Solución 9
\begin{minipage}{\textwidth}
\section*{Ejercicio 9}
\textbf{Ecuación:} \(1 - \cos^2 x - \cos x = 0\) en \([0,2\pi)\).

\textbf{Solución:}
\begin{enumerate}[label=\arabic*.]
    \item \textbf{Usar identidad:} \(1 - \cos^2 x = \sin^2 x\), pero conviene expresar todo en coseno.
    Sea \(u = \cos x\):
    \[
    1 - u^2 - u = 0 \quad\Longrightarrow\quad -u^2 - u + 1 = 0.
    \]
    Multiplicando por \(-1\):
    \[
    u^2 + u - 1 = 0.
    \]
    \item \textbf{Resolver la cuadrática:}
    \[
    u = \frac{-1 \pm \sqrt{1^2 - 4(1)(-1)}}{2} = \frac{-1 \pm \sqrt{5}}{2}.
    \]
    \item \textbf{Analizar los valores:}
    \begin{itemize}
        \item \(u = \frac{-1 - \sqrt{5}}{2} \approx -1.618\) (está fuera del rango de \(\cos x\)).
        \item \(u = \frac{-1 + \sqrt{5}}{2} = \frac{\sqrt{5} - 1}{2} \approx 0.618\) (válido).
    \end{itemize}
    Entonces \(\cos x = \frac{\sqrt{5} - 1}{2}\).
    \item \textbf{Resolver \(\cos x = a\) con \(a = \frac{\sqrt{5}-1}{2}\):}
    Usando la función \(\arccos\):
    \[
    x = \arccos\!\left(\frac{\sqrt{5} - 1}{2}\right) + 2k\pi,\qquad
    x = 2\pi - \arccos\!\left(\frac{\sqrt{5} - 1}{2}\right) + 2k\pi.
    \]
    En \([0,2\pi)\) con \(k=0\):
    \[
    x_1 = \arccos\!\left(\frac{\sqrt{5} - 1}{2}\right),\qquad
    x_2 = 2\pi - \arccos\!\left(\frac{\sqrt{5} - 1}{2}\right).
    \]
    \item \textbf{Conjunto solución:}
    \[
    \boxed{\displaystyle \arccos\!\left(\frac{\sqrt{5} - 1}{2}\right),\ 2\pi - \arccos\!\left(\frac{\sqrt{5} - 1}{2}\right)}
    \]
\end{enumerate}

\unitcircle{%
  ({1.2*cos(51.827)},{1.2*sin(51.827)})/\gamma,
  ({1.2*cos(308.173)},{1.2*sin(308.173)})/2\pi{-}\gamma%
} \\[2mm]
{\small $\gamma = \arccos\!\left(\frac{\sqrt{5} - 1}{2}\right)$.}
\end{minipage}

\vspace{1cm}

% Solución 10
\begin{minipage}{\textwidth}
\section*{Ejercicio 10}
\textbf{Ecuación:} \(2\sin^2 x + 5\sin x - 3 = 0\) en \([0,2\pi)\).

\textbf{Solución:}
\begin{enumerate}[label=\arabic*.]
    \item \textbf{Sustitución:} \(u = \sin x\) \(\Rightarrow\) \(2u^2 + 5u - 3 = 0\).
    \item \textbf{Factorización:} \((2u - 1)(u + 3) = 0\).
    \item \textbf{Raíces:} \(u = \frac{1}{2}\) o \(u = -3\).
    \item \textbf{Análisis:} \(\sin x = -3\) no tiene solución.\\
    \(\sin x = \frac{1}{2}\) es válida.
    \item \textbf{Resolver \(\sin x = \frac{1}{2}\):}
    \begin{align*}
        x &= \frac{\pi}{6} + 2k\pi,\\
        x &= \frac{5\pi}{6} + 2k\pi.
    \end{align*}
    En \([0,2\pi)\): \(x = \frac{\pi}{6},\ \frac{5\pi}{6}\).
    \item \textbf{Conjunto solución:}
    \[
    \boxed{\frac{\pi}{6},\ \frac{5\pi}{6}}
    \]
\end{enumerate}

\unitcircle{%
  ({1.2*cos(30)},{1.2*sin(30)})/\frac{\pi}{6},
  ({1.2*cos(150)},{1.2*sin(150)})/\frac{5\pi}{6}%
}
\end{minipage}

\vspace{1cm}

% Solución 11
\begin{minipage}{\textwidth}
\section*{Ejercicio 11}
\textbf{Ecuación:} \(\tan^2 x + 2\tan x - 3 = 0\) en \([0,2\pi)\).

\textbf{Solución:}
\begin{enumerate}[label=\arabic*.]
    \item \textbf{Sustitución:} \(u = \tan x\) \(\Rightarrow\) \(u^2 + 2u - 3 = 0\).
    \item \textbf{Factorización:} \((u + 3)(u - 1) = 0\).
    \item \textbf{Raíces:} \(u = -3\) o \(u = 1\).
    \item \textbf{Resolver \(\tan x = 1\):}
    Solución general: \(x = \frac{\pi}{4} + k\pi\). En \([0,2\pi)\): \(x = \frac{\pi}{4},\ \frac{5\pi}{4}\).
    \item \textbf{Resolver \(\tan x = -3\):}
    Usamos \(\arctan\):
    \[
    x = \arctan(-3) + k\pi.
    \]
    El valor principal \(\arctan(-3)\) es negativo (≈ \(-1.249\) rad). Para obtener ángulos en \([0,2\pi)\):
    \begin{itemize}
        \item Con \(k = 1\): \(x = \arctan(-3) + \pi = \pi - \arctan(3)\) (cuadrante II).
        \item Con \(k = 2\): \(x = \arctan(-3) + 2\pi = 2\pi - \arctan(3)\) (cuadrante IV).
    \end{itemize}
    \item \textbf{Conjunto solución:}
    \[
    \boxed{\frac{\pi}{4},\ \frac{5\pi}{4},\ \pi - \arctan(3),\ 2\pi - \arctan(3)}
    \]
\end{enumerate}

\unitcircle{%
  ({1.2*cos(45)},{1.2*sin(45)})/\frac{\pi}{4},
  ({1.2*cos(225)},{1.2*sin(225)})/\frac{5\pi}{4},
  ({1.2*cos(108.4349)},{1.2*sin(108.4349)})/\pi{-}\delta,
  ({1.2*cos(288.4349)},{1.2*sin(288.4349)})/2\pi{-}\delta%
} \\[2mm]
{\small $\delta = \arctan 3$.}
\end{minipage}

\vspace{1cm}

% Solución 12
\begin{minipage}{\textwidth}
\section*{Ejercicio 12}
\textbf{Ecuación:} \(\cos^2 x - \sin x = 0\) en \([0,2\pi)\).

\textbf{Solución:}
\begin{enumerate}[label=\arabic*.]
    \item \textbf{Usar identidad pitagórica:} \(\cos^2 x = 1 - \sin^2 x\).
    \[
    (1 - \sin^2 x) - \sin x = 0 \quad\Longrightarrow\quad -\sin^2 x - \sin x + 1 = 0.
    \]
    Multiplicando por \(-1\):
    \[
    \sin^2 x + \sin x - 1 = 0.
    \]
    \item \textbf{Sustitución:} \(u = \sin x\) \(\Rightarrow\) \(u^2 + u - 1 = 0\).
    \item \textbf{Resolver la cuadrática:}
    \[
    u = \frac{-1 \pm \sqrt{1 + 4}}{2} = \frac{-1 \pm \sqrt{5}}{2}.
    \]
    \item \textbf{Analizar valores:}
    \begin{itemize}
        \item \(u = \frac{-1 - \sqrt{5}}{2} \approx -1.618\) (fuera de rango).
        \item \(u = \frac{-1 + \sqrt{5}}{2} = \frac{\sqrt{5} - 1}{2} \approx 0.618\) (válido).
    \end{itemize}
    Entonces \(\sin x = \frac{\sqrt{5} - 1}{2}\).
    \item \textbf{Resolver \(\sin x = a\) con \(a = \frac{\sqrt{5}-1}{2}\):}
    \begin{align*}
        x &= \arcsin(a) + 2k\pi,\\
        x &= \pi - \arcsin(a) + 2k\pi.
    \end{align*}
    En \([0,2\pi)\) con \(k=0\):
    \[
    x_1 = \arcsin\!\left(\frac{\sqrt{5} - 1}{2}\right),\qquad
    x_2 = \pi - \arcsin\!\left(\frac{\sqrt{5} - 1}{2}\right).
    \]
    \item \textbf{Conjunto solución:}
    \[
    \boxed{\displaystyle \arcsin\!\left(\frac{\sqrt{5} - 1}{2}\right),\ \pi - \arcsin\!\left(\frac{\sqrt{5} - 1}{2}\right)}
    \]
\end{enumerate}

\unitcircle{%
  ({1.2*cos(38.1727)},{1.2*sin(38.1727)})/\varepsilon,
  ({1.2*cos(141.8273)},{1.2*sin(141.8273)})/\pi{-}\varepsilon%
} \\[2mm]
{\small $\varepsilon = \arcsin\!\left(\frac{\sqrt{5} - 1}{2}\right)$.}
\end{minipage}

\vspace{1cm}

% Solución 13
\begin{minipage}{\textwidth}
\section*{Ejercicio 13}
\textbf{Ecuación:} \(3\cos^2 x + 2\cos x - 1 = 0\) en \([0,2\pi)\).

\textbf{Solución:}
\begin{enumerate}[label=\arabic*.]
    \item \textbf{Sustitución:} \(u = \cos x\) \(\Rightarrow\) \(3u^2 + 2u - 1 = 0\).
    \item \textbf{Factorización:}
    \[
    3u^2 + 2u - 1 = (3u - 1)(u + 1) = 0.
    \]
    \item \textbf{Raíces:} \(u = \frac{1}{3}\) o \(u = -1\).
    \item \textbf{Resolver \(\cos x = -1\):}
    \[
    x = \pi + 2k\pi.
    \]
    En \([0,2\pi)\): \(x = \pi\).
    \item \textbf{Resolver \(\cos x = \frac{1}{3}\):}
    No es ángulo especial. Usamos \(\arccos\):
    \begin{align*}
        x &= \arccos\!\left(\frac{1}{3}\right) + 2k\pi,\\
        x &= 2\pi - \arccos\!\left(\frac{1}{3}\right) + 2k\pi.
    \end{align*}
    En \([0,2\pi)\) con \(k=0\):
    \[
    x_1 = \arccos\!\left(\frac{1}{3}\right),\qquad
    x_2 = 2\pi - \arccos\!\left(\frac{1}{3}\right).
    \]
    \item \textbf{Conjunto solución:}
    \[
    \boxed{\pi,\ \arccos\!\left(\frac{1}{3}\right),\ 2\pi - \arccos\!\left(\frac{1}{3}\right)}
    \]
\end{enumerate}

\unitcircle{%
  {(-1.2,0)/\pi},
  ({1.2*cos(70.5288)},{1.2*sin(70.5288)})/\zeta,
  ({1.2*cos(289.4712)},{1.2*sin(289.4712)})/2\pi{-}\zeta%
} \\[2mm]
{\small $\zeta = \arccos\frac13$.}
\end{minipage}

\vspace{1cm}

% Solución 14
\begin{minipage}{\textwidth}
\section*{Ejercicio 14}
\textbf{Ecuación:} \(2\sec^2 x + 5\sec x + 2 = 0\) en \([0,2\pi)\).

\textbf{Solución:}
\begin{enumerate}[label=\arabic*.]
    \item \textbf{Sustitución:} \(u = \sec x\) \(\Rightarrow\) \(2u^2 + 5u + 2 = 0\).
    \item \textbf{Factorización:} \((2u + 1)(u + 2) = 0\).
    \item \textbf{Raíces:} \(u = -\frac{1}{2}\) o \(u = -2\).
    \item \textbf{Volver a \(\cos x\):} \(\sec x = \frac{1}{\cos x}\).
    \begin{itemize}
        \item \(\sec x = -\frac{1}{2} \Rightarrow \cos x = -2\) (no válido).
        \item \(\sec x = -2 \Rightarrow \cos x = -\frac{1}{2}\).
    \end{itemize}
    \item \textbf{Resolver \(\cos x = -\frac{1}{2}\):}
    \begin{align*}
        x &= \frac{2\pi}{3} + 2k\pi,\\
        x &= \frac{4\pi}{3} + 2k\pi.
    \end{align*}
    En \([0,2\pi)\): \(x = \frac{2\pi}{3},\ \frac{4\pi}{3}\).
    \item \textbf{Conjunto solución:}
    \[
    \boxed{\frac{2\pi}{3},\ \frac{4\pi}{3}}
    \]
\end{enumerate}

\unitcircle{%
  ({1.2*cos(120)},{1.2*sin(120)})/\frac{2\pi}{3},
  ({1.2*cos(240)},{1.2*sin(240)})/\frac{4\pi}{3}%
}
\end{minipage}

\vspace{1cm}

% Solución 15
\begin{minipage}{\textwidth}
\section*{Ejercicio 15}
\textbf{Ecuación:} \(2\csc^2 x - 7\csc x + 3 = 0\) en \([0,2\pi)\).

\textbf{Solución:}
\begin{enumerate}[label=\arabic*.]
    \item \textbf{Sustitución:} \(u = \csc x\) \(\Rightarrow\) \(2u^2 - 7u + 3 = 0\).
    \item \textbf{Factorización:} \((2u - 1)(u - 3) = 0\).
    \item \textbf{Raíces:} \(u = \frac{1}{2}\) o \(u = 3\).
    \item \textbf{Volver a \(\sin x\):} \(\csc x = \frac{1}{\sin x}\).
    \begin{itemize}
        \item \(\csc x = \frac{1}{2} \Rightarrow \sin x = 2\) (no válido).
        \item \(\csc x = 3 \Rightarrow \sin x = \frac{1}{3}\).
    \end{itemize}
    \item \textbf{Resolver \(\sin x = \frac{1}{3}\):}
    \begin{align*}
        x &= \arcsin\!\left(\frac{1}{3}\right) + 2k\pi,\\
        x &= \pi - \arcsin\!\left(\frac{1}{3}\right) + 2k\pi.
    \end{align*}
    En \([0,2\pi)\) con \(k=0\):
    \[
    x_1 = \arcsin\!\left(\frac{1}{3}\right),\qquad
    x_2 = \pi - \arcsin\!\left(\frac{1}{3}\right).
    \]
    \item \textbf{Conjunto solución:}
    \[
    \boxed{\arcsin\!\left(\frac{1}{3}\right),\ \pi - \arcsin\!\left(\frac{1}{3}\right)}
    \]
\end{enumerate}

\unitcircle{%
  ({1.2*cos(19.4712)},{1.2*sin(19.4712)})/\eta,
  ({1.2*cos(160.5288)},{1.2*sin(160.5288)})/\pi{-}\eta%
} \\[2mm]
{\small $\eta = \arcsin\frac13$.}
\end{minipage}

\vspace{1cm}

% Solución 16
\begin{minipage}{\textwidth}
\section*{Ejercicio 16}
\textbf{Ecuación:} \(3\tan^2 x + \tan x - 2 = 0\) en \([0,2\pi)\).

\textbf{Solución:}
\begin{enumerate}[label=\arabic*.]
    \item \textbf{Sustitución:} \(u = \tan x\) \(\Rightarrow\) \(3u^2 + u - 2 = 0\).
    \item \textbf{Factorización:} \((3u - 2)(u + 1) = 0\).
    \item \textbf{Raíces:} \(u = \frac{2}{3}\) o \(u = -1\).
    \item \textbf{Resolver \(\tan x = -1\):}
    Solución general: \(x = -\frac{\pi}{4} + k\pi\). En \([0,2\pi)\):
    \begin{align*}
        k=1 &: x = -\frac{\pi}{4} + \pi = \frac{3\pi}{4},\\
        k=2 &: x = -\frac{\pi}{4} + 2\pi = \frac{7\pi}{4}.
    \end{align*}
    \item \textbf{Resolver \(\tan x = \frac{2}{3}\):}
    \begin{align*}
        x &= \arctan\!\left(\frac{2}{3}\right) + k\pi.
    \end{align*}
    En \([0,2\pi)\):
    \begin{align*}
        k=0 &: x = \arctan\!\left(\frac{2}{3}\right),\\
        k=1 &: x = \arctan\!\left(\frac{2}{3}\right) + \pi.
    \end{align*}
    \item \textbf{Conjunto solución:}
    \[
    \boxed{\frac{3\pi}{4},\ \frac{7\pi}{4},\ \arctan\!\left(\frac{2}{3}\right),\ \pi + \arctan\!\left(\frac{2}{3}\right)}
    \]
\end{enumerate}

\unitcircle{%
  ({1.2*cos(135)},{1.2*sin(135)})/\frac{3\pi}{4},
  ({1.2*cos(315)},{1.2*sin(315)})/\frac{7\pi}{4},
  ({1.2*cos(33.6901)},{1.2*sin(33.6901)})/\theta,
  ({1.2*cos(213.6901)},{1.2*sin(213.6901)})/\pi{+}\theta%
} \\[2mm]
{\small $\theta = \arctan\frac23$.}
\end{minipage}

\vspace{1cm}

% Solución 17
\begin{minipage}{\textwidth}
\section*{Ejercicio 17}
\textbf{Ecuación:} \(\cos^2 x - 2\cos x - 2 = 0\) en \([0,2\pi)\).

\textbf{Solución:}
\begin{enumerate}[label=\arabic*.]
    \item \textbf{Sustitución:} \(u = \cos x\) \(\Rightarrow\) \(u^2 - 2u - 2 = 0\).
    \item \textbf{Resolver la cuadrática:}
    \[
    u = \frac{2 \pm \sqrt{4 + 8}}{2} = \frac{2 \pm \sqrt{12}}{2} = \frac{2 \pm 2\sqrt{3}}{2} = 1 \pm \sqrt{3}.
    \]
    \item \textbf{Analizar valores:}
    \begin{itemize}
        \item \(1 + \sqrt{3} \approx 2.732\) (no válido).
        \item \(1 - \sqrt{3} \approx -0.732\) (válido).
    \end{itemize}
    Entonces \(\cos x = 1 - \sqrt{3}\).
    \item \textbf{Resolver \(\cos x = 1-\sqrt{3}\):}
    \begin{align*}
        x &= \arccos(1 - \sqrt{3}) + 2k\pi,\\
        x &= 2\pi - \arccos(1 - \sqrt{3}) + 2k\pi.
    \end{align*}
    En \([0,2\pi)\) con \(k=0\):
    \[
    x_1 = \arccos(1 - \sqrt{3}),\qquad
    x_2 = 2\pi - \arccos(1 - \sqrt{3}).
    \]
    \item \textbf{Conjunto solución:}
    \[
    \boxed{\arccos(1 - \sqrt{3}),\ 2\pi - \arccos(1 - \sqrt{3})}
    \]
\end{enumerate}

\unitcircle{%
  ({1.2*cos(137.304)},{1.2*sin(137.304)})/\iota,
  ({1.2*cos(222.696)},{1.2*sin(222.696)})/2\pi{-}\iota%
} \\[2mm]
{\small $\iota = \arccos(1 - \sqrt{3})$.}
\end{minipage}

\vspace{1cm}

% Solución 18
\begin{minipage}{\textwidth}
\section*{Ejercicio 18}
\textbf{Ecuación:} \(5\sin^2 x - 4\sin x - 1 = 0\) en \([0,2\pi)\).

\textbf{Solución:}
\begin{enumerate}[label=\arabic*.]
    \item \textbf{Sustitución:} \(u = \sin x\) \(\Rightarrow\) \(5u^2 - 4u - 1 = 0\).
    \item \textbf{Factorización:} \((5u + 1)(u - 1) = 0\).
    \item \textbf{Raíces:} \(u = -\frac{1}{5}\) o \(u = 1\).
    \item \textbf{Resolver \(\sin x = 1\):}
    \[
    x = \frac{\pi}{2} + 2k\pi.
    \]
    En \([0,2\pi)\): \(x = \frac{\pi}{2}\).
    \item \textbf{Resolver \(\sin x = -\frac{1}{5}\):}
    Usamos \(\arcsin\):
    \begin{align*}
        x &= \arcsin\!\left(-\frac{1}{5}\right) + 2k\pi,\\
        x &= \pi - \arcsin\!\left(-\frac{1}{5}\right) + 2k\pi.
    \end{align*}
    Como \(\arcsin(-a) = -\arcsin(a)\), las expresiones quedan:
    \begin{align*}
        x &= -\arcsin\!\left(\frac{1}{5}\right) + 2k\pi,\\
        x &= \pi + \arcsin\!\left(\frac{1}{5}\right) + 2k\pi.
    \end{align*}
    Para estar en \([0,2\pi)\):
    \begin{itemize}
        \item De la primera, con \(k=1\): \(x = 2\pi - \arcsin\!\left(\frac{1}{5}\right)\).
        \item De la segunda, con \(k=0\): \(x = \pi + \arcsin\!\left(\frac{1}{5}\right)\).
    \end{itemize}
    \item \textbf{Conjunto solución:}
    \[
    \boxed{\frac{\pi}{2},\ \pi + \arcsin\!\left(\frac{1}{5}\right),\ 2\pi - \arcsin\!\left(\frac{1}{5}\right)}
    \]
\end{enumerate}

\unitcircle{%
  {(0,1.2)/\frac{\pi}{2}},
  ({1.2*cos(191.537)},{1.2*sin(191.537)})/\pi{+}\kappa,
  ({1.2*cos(348.463)},{1.2*sin(348.463)})/2\pi{-}\kappa%
} \\[2mm]
{\small $\kappa = \arcsin\frac15$.}
\end{minipage}

\vspace{1cm}

% Solución 19
\begin{minipage}{\textwidth}
\section*{Ejercicio 19}
\textbf{Ecuación:} \(3\tan^2 x - 1 = 0\) en \([0,2\pi)\).

\textbf{Solución:}
\begin{enumerate}[label=\arabic*.]
    \item \textbf{Despejar \(\tan^2 x\):}
    \[
    3\tan^2 x = 1 \quad\Longrightarrow\quad \tan^2 x = \frac{1}{3}.
    \]
    \item \textbf{Tomar raíz cuadrada:}
    \[
    \tan x = \pm \frac{1}{\sqrt{3}} = \pm \frac{\sqrt{3}}{3}.
    \]
    \item \textbf{Resolver \(\tan x = \frac{\sqrt{3}}{3}\):}
    Ángulo de referencia \(\frac{\pi}{6}\). Tangente positiva en I y III:
    \begin{align*}
        x &= \frac{\pi}{6} + k\pi,\\
        \text{en }[0,2\pi): &\quad \frac{\pi}{6},\ \frac{7\pi}{6}.
    \end{align*}
    \item \textbf{Resolver \(\tan x = -\frac{\sqrt{3}}{3}\):}
    Ángulo de referencia \(\frac{\pi}{6}\). Tangente negativa en II y IV:
    \begin{align*}
        x &= -\frac{\pi}{6} + k\pi,\\
        \text{en }[0,2\pi): &\quad \frac{5\pi}{6} \ (k=1),\ \frac{11\pi}{6} \ (k=2).
    \end{align*}
    \item \textbf{Conjunto solución:}
    \[
    \boxed{\frac{\pi}{6},\ \frac{5\pi}{6},\ \frac{7\pi}{6},\ \frac{11\pi}{6}}
    \]
\end{enumerate}

\unitcircle{%
  ({1.2*cos(30)},{1.2*sin(30)})/\frac{\pi}{6},
  ({1.2*cos(150)},{1.2*sin(150)})/\frac{5\pi}{6},
  ({1.2*cos(210)},{1.2*sin(210)})/\frac{7\pi}{6},
  ({1.2*cos(330)},{1.2*sin(330)})/\frac{11\pi}{6}%
}
\end{minipage}

\vspace{1cm}

% Solución 20
\begin{minipage}{\textwidth}
\section*{Ejercicio 20}
\textbf{Ecuación:} \(2\sin^2 x - \sin x - 1 = 0\) en \([0,2\pi)\).

\textbf{Solución:}
\begin{enumerate}[label=\arabic*.]
    \item \textbf{Sustitución:} \(u = \sin x\) \(\Rightarrow\) \(2u^2 - u - 1 = 0\).
    \item \textbf{Factorización:} \((2u + 1)(u - 1) = 0\).
    \item \textbf{Raíces:} \(u = -\frac{1}{2}\) o \(u = 1\).
    \item \textbf{Resolver \(\sin x = 1\):}
    \[
    x = \frac{\pi}{2} + 2k\pi.
    \]
    En \([0,2\pi)\): \(x = \frac{\pi}{2}\).
    \item \textbf{Resolver \(\sin x = -\frac{1}{2}\):}
    Ángulo de referencia \(\frac{\pi}{6}\). Seno negativo en III y IV:
    \begin{align*}
        x &= \pi + \frac{\pi}{6} + 2k\pi = \frac{7\pi}{6} + 2k\pi,\\
        x &= 2\pi - \frac{\pi}{6} + 2k\pi = \frac{11\pi}{6} + 2k\pi.
    \end{align*}
    En \([0,2\pi)\): \(x = \frac{7\pi}{6},\ \frac{11\pi}{6}\).
    \item \textbf{Conjunto solución:}
    \[
    \boxed{\frac{\pi}{2},\ \frac{7\pi}{6},\ \frac{11\pi}{6}}
    \]
\end{enumerate}

\unitcircle{%
  {(0,1.2)/\frac{\pi}{2}},
  ({1.2*cos(210)},{1.2*sin(210)})/\frac{7\pi}{6},
  ({1.2*cos(330)},{1.2*sin(330)})/\frac{11\pi}{6}%
}
\end{minipage}

\end{document}